\subsection{\texorpdfstring{Bounded subsets of \(\mathcal{L}(\mathbf{E};\mathbf{F})\) (quasi-complete case)}{Bounded subsets of \textbackslash mathcal\{L\}(\textbackslash mathbf\{E\};\textbackslash mathbf\{F\}) (quasi-complete case)}}

THEOREM 2. --- Let E be a locally convex Hausdorff space, F a locally convex space and S a family of closed, convex, balanced, bounded and semi-complete subsets of E (III, p.~7). Every simply bounded subset H of \(\mathcal{L}(\mathrm{E};\mathrm{F})\) is bounded for the S-topology.

Let \(A \in S\) . The space \(E_A\) is then a Banach space (III, p.~8, corollary), hence barrelled. On the other hand, the canonical image of H in \(\mathcal{L}(E_A; F)\) is simply bounded, hence equicontinuous (III, p.~25, th. 1). Consequently, the set of all \(u(x)\) for \(u \in H\) and \(x \in A\) , is bounded in F, which proves that H is bounded for the \(S\) -topology.

COROLLARY 1. --- Let E be a locally convex Hausdorff space, F a locally convex space, and S a family of bounded subsets of E. If E is semi-complete, then every simply bounded subset of \(\mathcal{L}(\mathrm{E};\mathrm{F})\) is bounded for the S-topology.

It is enough to apply th. 2, after replacing the sets of S by their closed, convex, balanced envelopes, since this does not change the S-topology.

When E is semi-complete (for example quasi-complete), we can talk of the bounded subsets of \(\mathcal{L}(\mathrm{E};\mathrm{F})\) without specifying the S-topology, since these are the same for all the S-topologies whenever S is a cover of E.

COROLLARY 2. --- Every semi-complete bornological space is barrelled.

Every simply bounded subset of the dual of such a space is strongly bounded (cor. 1), hence equicontinuous (III, p.~22, prop. 10).

COROLLARY 3. --- Let E be a locally convex space. Every subset of E which is bounded for \(\sigma(E, E')\) is bounded.

Let A be a subset of E. Saying that A is bounded for \(\sigma(\mathrm{E}, \mathrm{E}')\) means that every continuous linear form on E is bounded on A; Saying that A is bounded means that every continuous semi-norm on E is bounded on A. Let N be the closure of 0 in E and \(\pi\) the canonical mapping from E onto E/N. The continuous linear forms on E are the mappings of the form \(f \circ \pi\) with \(f \in (\mathrm{E}/\mathrm{N})'\) and we have an analogous characterization of continuous semi-norms on E. Replacing E by E/N and A by \(\pi(\mathrm{A})\) we can thus limit ourselves to the case where E is Hausdorff.

Let \(\mathfrak{S}\) be the set of equicontinuous subsets of \(E'\) ; when \(E'\) is assigned the topology \(\sigma(E', E)\) , \(E\) can be identified with \((E')_{\mathfrak{S}}'\) (III, p.~19, cor. 1). Every closed equicontinuous subset of \(E'\) is compact for \(\sigma(E', E)\) (III, p.~17, cor. 2), hence complete for \(\sigma(E', E)\) . It is now enough to apply th. 2.

